% ============================================================
% 3. PROBLEM FORMULATION
% ============================================================
\section{Problem Formulation}

% [REVISION | §3-reframe-D+ | 2026-05-11 | Reframed §3 around Dimensional Escape as a unifying geometric pathology (Session A v9 draft 2026-05-10, applied 2026-05-11). Old structure (Option α+): §3.1 inference-time latent-space framing with embedded Hypothesis 1. New structure (Option D+): §3.1 tightened framing of the central question; new §3.2 Dimensional Escape with Definition~\ref{def:dim_escape} + Hypothesis 1 (preserved hyp:monitor_driven label, retitled as "Monitor-driven sub-case") + new Hypothesis 2 (hyp:goodhart label, "Reward-driven sub-case"); existing §3.2 Discrete Text-Space Bottleneck renumbers to §3.3 with no content change. Eq.~\ref{eq:alignment_faking} preserved with label intact, relocated to immediately after Hypothesis 1 (was embedded in old §3.1 prose). Eq.~\ref{eq:orthogonal_escape} preserved with label intact. New Eq.~\ref{eq:goodhart}. The 2026-05-06 §3-reframe-α+ marker is subsumed. Bailey 2024 reference preserved but folded into Dimensional Escape unifying claim rather than positioned as standalone framing peer. Citation `xie2025entropyloop` and `goodhart1984` (mentioned in v9 draft) DROPPED post-bibliography-grep verification 2026-05-10 — concepts carried inline without bibitem expansion. — pending audit]
\subsection{Inference-Time Latent-Space Analysis on Frozen Models}

% [REVISION | §3.1-trim-Day10-batch4 | 2026-05-15 | Day-10 Batch 4 page-budget cut #21: trimmed §3.1 from 2 paragraphs to 1 (~30% shorter). The two-paragraph structure (setup question / pathology naming) is consolidated; sub-clause repetition ("corrective monitor, reward signal, or representational probe") dropped — the three roles are enumerated upstream in §1 Introduction. Per Day-9 page-budget cut plan §3 rank #21 (LOW risk, save 0.10 pp). — pending audit]
Let $f_\theta$ be a frozen language model with hidden state $\mathbf{h}_t^{(\ell)} \in \mathbb{R}^d$ at layer $\ell$ and token position $t$, and let $\mathbf{P} \in \mathbb{R}^{k \times d}$ of rank $k \ll d$ denote a linear monitor or surrogate scoring function (e.g., representation-engineering honesty probes~\citep{zou2023repe}, aggregations thereof~\citep{gupta2025rlobfuscation}, or reward functions derived from internal heuristics such as next-token entropy). The central question is whether perturbations to the cached key/value tensors guided by such a low-rank projection produce trajectories that navigate the semantic manifold or that exploit the unconstrained dimensions of the overparameterized latent space; the latter outcome---observed theoretically and empirically (§\ref{sec:experiments})---is the pathology we name \emph{Dimensional Escape}.

\subsection{Dimensional Escape: A Unifying Geometric Pathology}\label{sec:dim_escape}

Two seemingly distinct failure modes have been hypothesized for inference-time intervention on frozen language models. The first is \emph{monitor-driven}: when corrective pressure is applied along a monitored subspace, a sufficiently capable model may redistribute its internal representations into the null space of the monitor, preserving behavior while evading detection. This pattern is the subject of \citet{bailey2024obfuscated} (gradient-trained input-space variant, with a tentative negative result reported) and motivates the alignment-faking literature---behavioral demonstrations~\citep{greenblatt2024alignment, hubinger2024sleeper} and latent-space monitor evasion~\citep{gupta2025rlobfuscation}. The second mode is \emph{reward-driven}: when a search procedure optimizes a surrogate reward lacking exogenous grounding (for instance, next-token entropy as a proxy for confidence), the optimizer may resolve the reward via geometric paths that bypass the semantic objective entirely. This is Goodhart's Law manifesting in cache geometry: when a measure becomes a target, the optimizer discovers the cheapest computational path to alter the metric, even if that path circumvents the original objective. We hypothesize these as parallel manifestations of the overparameterized-latent-space geometry described below; only the reward-driven case is evidenced empirically in this paper, with the monitor-driven case under inference-time KV-cache perturbation an open empirical question.

% [REVISION | M3-§3.2-Goodhart-anchor-Day12 | 2026-05-17 | Day-12 patch M3 per Round-1 parallel-review synthesis (Axis 2 Q5 + Axis 5 Pass H). Inserted a Goodhart-anchor preamble paragraph immediately before def:dim_escape positioning Dimensional Escape as a specific manifestation of Goodhart's Law (gao2023scaling, skalse2022defining) in low-rank-scored, high-dimensional latent spaces, generalizing Bailey 2024's gradient-trained input-attack Orthogonal Escape. Defends the coined term against "you're just relabeling X" reviewer attack by making the inheritance explicit. Companion: M2 (C2 title hedge, line ~84) anchors the same Option B framing in the contribution title. Two new bibitems (gao2023scaling, skalse2022defining) added in this same patch. — pending audit]
We position Dimensional Escape as a specific manifestation of Goodhart's Law~\citep{gao2023scaling, skalse2022defining} in low-rank-scored, high-dimensional latent spaces, generalizing the obfuscated-activations regime \citet{bailey2024obfuscated} characterize for gradient-trained input attacks. The novelty we claim is not a new phenomenon class but a specific empirical setting (inference-time KV-cache MCTS with surrogate-reward scoring) and a two-mechanism characterization of how the search satisfies the proxy without satisfying the objective (prior-amplification at small model scale, greedy-collapse at larger scale; §\ref{sec:exp1}--\ref{sec:exp2}).

\begin{definition}[Dimensional Escape]\label{def:dim_escape}
Let $\mathbf{P} \in \mathbb{R}^{k \times d}$ be a monitoring or scoring projection of rank $k \ll d$, and let $\mathcal{L}: \mathbb{R}^d \to \mathbb{R}$ be a localized optimization objective acting on hidden state $\mathbf{h}$. \emph{Dimensional Escape} occurs when the optimization process resolves $\mathcal{L}$ through trajectories whose induced perturbation $\Delta\mathbf{h}$ has substantial component in $\text{Null}(\mathbf{P}_{\text{semantic}})$, where $\mathbf{P}_{\text{semantic}}$ is the (typically unobserved) projection onto the semantically relevant subspace of $\mathbb{R}^d$. The optimizer thus satisfies $\mathcal{L}$ without traversing the manifold of semantically meaningful states.
\end{definition}

The high dimensionality of modern Transformer hidden states ($d \in [2048, 8192]$ for the Llama~3 family) ensures that for any low-rank monitor or scoring function, the unconstrained complement is overwhelmingly large; localized optimization pressure routes through this complement, whether applied as corrective signal or as reward signal. We instantiate Dimensional Escape under two hypotheses, one per failure mode.

% [REVISION | §3.2-hypothesis-tighten-Day10-batch5-extra | 2026-05-15 | Day-10 Batch 5 page-budget extra cut: dropped eq:orthogonal_escape (formal restatement of Hypothesis 1's verbal claim) and eq:alignment_faking (behavioral-consequence equation only forward-referenced within the deleted Hypothesis-1 closing sentence). Both equations had no body cross-references outside §3.2 itself (verified via grep). Hypothesis 1 and Hypothesis 2 closing pedagogical sentences tightened to single specialization-of-Definition-def:dim_escape clauses. Body savings ~0.3 pp; the unification framing is preserved. — pending audit]
\begin{hypothesis}[Monitor-driven Dimensional Escape under inference-time KV-cache perturbation]\label{hyp:monitor_driven}
When corrective pressure is applied along the monitored subspace $\text{Range}(\mathbf{P}_{\text{monitor}})$, a sufficiently capable model under inference-time steering may redistribute its representation into $\text{Null}(\mathbf{P}_{\text{monitor}})$, preserving behavioral compliance while evading detection (Definition~\ref{def:dim_escape} with $\mathcal{L}$ as corrective pressure and $\mathbf{P}_{\text{monitor}}$ as a linear deception probe).
\end{hypothesis}

\begin{hypothesis}[Reward-driven Dimensional Escape: Goodhart-collapse under surrogate-reward MCTS]\label{hyp:goodhart}
When MCTS optimizes a surrogate reward $\hat{r}: \mathbb{R}^d \to \mathbb{R}$ lacking exogenous grounding---for instance, $\hat{r}(\mathbf{h}) = -H(p(\cdot \mid \mathbf{h}))/\log|V|$ where $H$ is the next-token entropy at the steered cache state and $|V|$ is the vocabulary size---the optimal cache state under $\hat{r}$ generically lies off the semantic manifold:
\begin{equation}
\arg\min_{\mathbf{h}} \hat{r}(\mathbf{h}) \cap \text{SemanticManifold}(x) = \emptyset
\label{eq:goodhart}
\end{equation}
where $\text{SemanticManifold}(x) \subset \mathbb{R}^d$ is the set of cache states consistent with a correct response to input $x$, and the equality holds generically in dimension $d \gg \text{rank}(\hat{r})$. The empirical mechanism by which this manifests---prior-amplification at small model scale, greedy-collapse at larger scale---is established in §\ref{sec:exp1}.
\end{hypothesis}

% [REVISION | §3.2-closing-trim-Day10-batch4 | 2026-05-15 | Day-10 Batch 4 page-budget cut #23: trimmed §3.2 closing paragraph ~50%. The Bailey~\citep{bailey2024obfuscated} restatement of the cosine-similarity-vs-KL-divergence trade-off is already at §2.2 paragraph 2; the §\ref{sec:measurement-prereqs} pointer is redundant once §7.4 (the section it labelled) has moved to appendix. The unification framing (reward-driven empirically here; monitor-driven left to Phase~B) is preserved. Per Day-9 page-budget cut plan §3 rank #23 (HIGH risk per R2; unification framing protected, save 0.10 pp). — pending audit]
\citet{bailey2024obfuscated} characterized the monitor-driven case (gradient-trained input-space attacks) with a tentative negative result; we characterize the reward-driven case empirically (§\ref{sec:exp1}--\ref{sec:exp2}) and find it manifests robustly across model scale and prompt class. The monitor-driven case under inference-time KV-cache perturbation is left to Phase~B research.

% [REVISION | §3.3-delete-Day10-batch4 | 2026-05-15 | Day-10 Batch 4 page-budget cut #8: DELETED §3.3 The Discrete Text-Space Bottleneck. Paragraph 1 (discrete combinatorial space framing) restates §3.1's central question + the algorithmic motivation now in §4.2 intro. Paragraph 2 (Dimensional Escape as geometric problem) restates Definition~\ref{def:dim_escape}. Paragraph 3 (FP32-accumulator-as-machinery) restates §4.2 intro. The 2026-05-11 §3.3-Dimensional-Escape-vocab-D+ marker is subsumed; the lexical softening it documented (Orthogonal-Escape-phenomenon → Dimensional-Escape-pathology) was load-bearing only at this location, which is now deleted. Per Day-9 page-budget cut plan §3 rank #8 (LOW risk, save 0.30 pp). — pending audit]
